\documentclass[12pt,a4paper]{article}

% 中文支持 fontset=none (CJK 字体族手动设置; AutoFakeBold 在 \setCJKmainfont 选项中启用)
\usepackage[fontset=none,UTF8]{ctex}
% 手动设置 CJK 字体族 (fontset=none 时必须)
\setCJKmainfont[AutoFakeBold,AutoFakeSlant]{SimSun}
\setCJKsansfont[AutoFakeBold,AutoFakeSlant]{SimHei}
\setCJKfamilyfont{zhsong}[AutoFakeBold]{SimSun}
\setCJKfamilyfont{zhhei}[AutoFakeBold]{SimHei}
\setCJKfamilyfont{zhkai}[AutoFakeBold]{KaiTi}
\setCJKfamilyfont{zhfs}[AutoFakeBold]{FangSong}

% 微调排版 + 数学增强
\usepackage[protrusion=true,expansion=false]{microtype}
% 数学包
\usepackage{amsmath,amssymb,amsfonts}
\usepackage{mathtools}

% 算法 (algorithm2e 替代旧 algorithm+algorithmic 组合)
\usepackage[ruled,linesnumbered]{algorithm2e}

% 图表
\usepackage{graphicx}
\usepackage{float}
\usepackage{subcaption}
\usepackage{booktabs}
\usepackage{multirow}

% 引用和链接
\usepackage[numbers,sort&compress]{natbib}
\usepackage[colorlinks=true,linkcolor=blue,citecolor=blue,urlcolor=blue]{hyperref}

% 页面设置
\usepackage[margin=2.5cm]{geometry}

% 摘要左右边距与正文统一：去掉 article.cls 默认 \quotation 造成的左右内缩（约 2.5em），
% 使摘要正文占满全文宽度，与正文段落左右边距一致。
\renewenvironment{abstract}{%
  \small
  \begin{center}%
    {\bfseries\abstractname\vspace{-.5em}}%
  \end{center}%
  \par\noindent\ignorespaces
}{%
  \par
}

% 代码
\usepackage{listings}
\usepackage{xcolor}

\lstset{
    language=Python,
    basicstyle=\small\ttfamily,
    keywordstyle=\color{blue},
    commentstyle=\color{gray},
    stringstyle=\color{red},
    numbers=left,
    numberstyle=\tiny\color{gray},
    frame=single,
    breaklines=true,
    captionpos=b
}

% 标题信息
\title{基于[方法名称]的[问题名称]研究}
\author{作者姓名\\[0.5em]
\small 单位名称\\
\small \texttt{email@example.com}}
\date{\today}

% 定理类环境（论文常用）
\newtheorem{definition}{Definition}
\newtheorem{theorem}{Theorem}
\newtheorem{lemma}{Lemma}
\newtheorem{corollary}{Corollary}
\newtheorem{proposition}{Proposition}
\newtheorem{example}{Example}
\newtheorem{remark}{Remark}

\begin{document}

\maketitle

\begin{abstract}
本文研究了[问题描述]。针对现有方法的[局限性]，提出了一种基于[方法]的[改进方法]。该方法的主要特点包括：(1) [特点1]；(2) [特点2]；(3) [特点3]。在[数据集/实验环境]上的实验结果表明，所提方法在[指标1]上达到了[数值]，相比基线方法提升了[百分比]。理论分析和实验验证均表明，该方法具有[优势]。
\textbf{关键词}：关键词1；关键词2；关键词3；关键词4
\end{abstract}

\section{引言}

\subsection{研究背景}
[描述研究问题的背景和重要性]

\subsection{问题描述}
[正式描述要解决的问题]

\subsection{主要贡献}
本文的主要贡献包括：
\begin{itemize}
    \item 贡献1：[详细描述]
    \item 贡献2：[详细描述]
    \item 贡献3：[详细描述]
\end{itemize}

\subsection{论文组织}
本文其余部分组织如下：第2节回顾相关工作；第3节描述所提方法；第4节给出实验结果；第5节得出结论。

\section{相关工作}

\subsection{[相关领域1]}
[描述相关工作1]

\subsection{[相关领域2]}
[描述相关工作2]

\section{方法}

\subsection{问题定义}
\begin{definition}[问题名称]
给定输入 $X = \{x_1, x_2, \ldots, x_n\}$，目标是找到映射 $f: X \rightarrow Y$，使得...
\end{definition}

\subsection{模型架构}
[描述整体模型架构]

\subsection{目标函数}
优化目标为：
\begin{equation}
\mathcal{L} = \frac{1}{N} \sum_{i=1}^{N} \ell(f(x_i), y_i) + \lambda \Omega(\theta)
\label{eq:loss}
\end{equation}
其中，$\ell$ 是损失函数，$\Omega$ 是正则化项，$\lambda$ 是平衡参数。

\subsection{算法流程}
\begin{algorithm}[H]
\caption{算法名称}
\label{alg:algorithm}
\KwIn{训练数据 $D = \{(x_i, y_i)\}_{i=1}^N$，学习率 $\eta$，迭代次数 $T$}
\KwOut{模型参数 $\theta^*$}
初始化参数 $\theta$\;
\For{$t = 1$ \KwTo $T$}{
    采样mini-batch $B \subset D$\;
    计算梯度 $g_t = \nabla_\theta \mathcal{L}_B(\theta)$\;
    更新参数 $\theta \leftarrow \theta - \eta g_t$\;
}
\Return $\theta^* = \theta$\;
\end{algorithm}

\section{实验}

\subsection{实验设置}
\subsubsection{数据集}
[描述使用的数据集]

\subsubsection{评估指标}
[描述评估指标]

\subsubsection{基线方法}
[描述基线方法]

\subsubsection{实现细节}
[描述实现细节]

\subsection{实验结果}

\begin{table}[H]
\centering
\caption{主要实验结果}
\label{tab:main_results}
\begin{tabular}{lcc}
\toprule
方法 & 指标1 & 指标2 \\
\midrule
基线1 & xx.x & xx.x \\
基线2 & xx.x & xx.x \\
基线3 & xx.x & xx.x \\
\textbf{本文方法} & \textbf{xx.x} & \textbf{xx.x} \\
\bottomrule
\end{tabular}
\end{table}

\subsection{分析与讨论}

\subsubsection{消融实验}
[描述消融实验结果]

\subsubsection{参数敏感性分析}
[描述参数敏感性分析]

\subsubsection{案例分析}
[描述案例分析]

\section{结论}
本文提出了[方法名称]，用于解决[问题]。实验结果表明，所提方法在[指标]上取得了优异的性能。未来工作将包括：(1) [未来工作1]；(2) [未来工作2]。

\section*{致谢}
[致谢内容]

\bibliographystyle{plainnat}
\bibliography{references}

\end{document}